\documentclass[11pt,aspectratio=169,table]{beamer}

\usetheme{Madrid}
\usecolortheme{beaver}
\setbeamertemplate{navigation symbols}{}

\usepackage[utf8]{inputenc}
\usepackage[T1]{fontenc}
\usepackage{amsmath,amssymb}
\usepackage{booktabs}
\usepackage{listings}
\usepackage{xcolor}
\usepackage{tcolorbox}
\tcbuselibrary{skins,breakable}

\definecolor{codebg}{RGB}{245,245,245}
\definecolor{codekw}{RGB}{0,0,180}
\lstdefinestyle{latexcode}{
  backgroundcolor=\color{codebg},
  basicstyle=\ttfamily\scriptsize,
  keywordstyle=\color{codekw}\bfseries,
  commentstyle=\color{gray},
  frame=single,
  rulecolor=\color{gray!40},
  breaklines=true,
  showstringspaces=false,
  xleftmargin=4pt,
  xrightmargin=4pt,
}
\lstset{style=latexcode}

\newtcolorbox{outputbox}{
  colback=white, colframe=black!65, boxrule=0.5pt, arc=1pt,
  left=8pt, right=8pt, top=5pt, bottom=4pt,
  title=\textsc{Compiled Output}, fonttitle=\bfseries\scriptsize,
  coltitle=white, colbacktitle=black!65
}
\newcommand{\whybox}[1]{%
  \vspace{2pt}
  {\scriptsize\raggedright\textcolor{red!55!black}{\textbf{Why \& What:}}~#1\par}
}

\newtcolorbox{exercisebox}{
  colback=orange!8, colframe=orange!70!black, boxrule=0.6pt,
  arc=2pt, left=6pt, right=6pt, top=4pt, bottom=4pt,
  title={Hands-on Exercise}, fonttitle=\bfseries, coltitle=white,
  colbacktitle=orange!70!black
}
\newtcolorbox{homeworkbox}{
  colback=blue!6, colframe=blue!50!black, boxrule=0.6pt,
  arc=2pt, left=6pt, right=6pt, top=4pt, bottom=4pt,
  title={Homework}, fonttitle=\bfseries, coltitle=white,
  colbacktitle=blue!50!black
}

\title[\texttt{\textbackslash begin} to \texttt{\textbackslash end}]{From \texttt{\textbackslash begin} to \texttt{\textbackslash end}}
\subtitle{The Complete LaTeX Guide \\[2pt] \normalsize Day 6 of 7: Journal Templates \& Professional Formatting}
\author[A.\ Rahman]{Atikur Rahman}
\institute[White Globe Initiative]{White Globe Initiative}

\setbeamertemplate{footline}{
  \leavevmode
  \hbox{%
  \begin{beamercolorbox}[wd=.38\paperwidth,ht=2.5ex,dp=1.125ex,leftskip=1em]{author in head/foot}%
    \usebeamerfont{author in head/foot}White Globe Initiative
  \end{beamercolorbox}%
  \begin{beamercolorbox}[wd=.42\paperwidth,ht=2.5ex,dp=1.125ex,center]{title in head/foot}%
    \usebeamerfont{title in head/foot}From \texttt{\textbackslash begin} to \texttt{\textbackslash end} --- Day 6
  \end{beamercolorbox}%
  \begin{beamercolorbox}[wd=.20\paperwidth,ht=2.5ex,dp=1.125ex,right,rightskip=1em]{date in head/foot}%
    \usebeamerfont{date in head/foot}\insertframenumber{} / \inserttotalframenumber
  \end{beamercolorbox}}%
  \vskip0pt
}

\begin{document}

\begin{frame}[plain]
  \begin{center}
    \vspace{4pt}
    {\Huge\bfseries\usebeamercolor[fg]{title}From \texttt{\textbackslash begin} to \texttt{\textbackslash end}}\\[6pt]
    {\Large\itshape The Complete LaTeX Guide}\\[18pt]
    {\large\bfseries Day 6 of 7 --- Journal Templates \& Professional Formatting}\\[22pt]
    \rule{0.7\textwidth}{0.4pt}\\[14pt]
    \begin{tabular}{rl}
      \textbf{Instructor:}       & Atikur Rahman \\[3pt]
      \textbf{Organized by:}     & White Globe Initiative \\[3pt]
      \textbf{Instructor Email:} & \texttt{atik@whiteglobeinitiative.org} \\[3pt]
      \textbf{General Contact:}  & \texttt{contact@whiteglobeinitiative.org} \\[3pt]
      \textbf{Format:}           & 7-Day Intensive Course \\[3pt]
      \textbf{Session:}          & Day 6 of 7 \\
    \end{tabular}\\[18pt]
    \rule{0.7\textwidth}{0.4pt}
  \end{center}
\end{frame}

\begin{frame}{Today's Agenda}
  \tableofcontents
\end{frame}

\section{Recap \& Today's Goal}

\begin{frame}{Quick Recap: Day 5}
  \small
  Yesterday: table of contents, footnotes, \texttt{hyperref} for
  clickable links, splitting a long paper across files with
  \texttt{\textbackslash input}, and page numbering/headers.
  \vspace{6pt}

  \textbf{Today's goal:} you ultimately need to submit to a specific
  journal or working-paper series. Leave today fluent in
  \textbf{adapting an existing template} rather than building
  formatting from scratch.
\end{frame}

\section{The Overleaf Template Gallery}

\begin{frame}{Browsing the Template Gallery}
  \small
  \texttt{overleaf.com/latex/templates} -- filter by \textbf{``Academic
  Journal''} to see hundreds of ready-made options. Economics-relevant
  choices: \textbf{AEA-style templates} for American Economic
  Association journals (AER, AEJ series); \textbf{Elsevier's
  \texttt{elsarticle} class}, used by many economics and finance
  journals (e.g.\ \emph{Journal of Development Economics}); and
  generic working-paper templates for institutions like NBER, IMF, or
  your own university series. Every template opens directly in
  Overleaf, pre-loaded with the correct class, packages, and
  bibliography style -- you never assemble one from scratch.
\end{frame}

\section{Adapting a Template}

\begin{frame}[fragile]{Filling In the Template's Metadata}
  \begin{lstlisting}
\title{Remittances and Consumption Smoothing}
\author{Jane Economist}
\begin{abstract}
This paper examines remittance inflows...
\end{abstract}
\JEL{F24, D14, O15}
\keywords{remittances, migration}
  \end{lstlisting}
  \begin{outputbox}
    \scriptsize
    \textbf{Remittances and Consumption Smoothing} --- Jane Economist\\[3pt]
    \textit{Abstract.} This paper examines remittance inflows...\\[3pt]
    \textit{JEL:} F24, D14, O15 \quad \textit{Keywords:} remittances, migration
  \end{outputbox}
  \whybox{Journal templates add fields beyond plain \texttt{article}:
  \textbf{abstract}, \textbf{JEL codes}, \textbf{keywords} -- rendered
  automatically as shown above. Command names vary by template --
  always check its own comments first. JEL codes are
  economics-specific and expected by nearly every journal.}
\end{frame}

\section{Page Size, Margins \& Colour}

\begin{frame}[fragile]{When a Template Doesn't Match House Style}
  \begin{lstlisting}
\rowcolor{gray!10}
Bangladesh & 6.0\% & 5.4\% \\
India & 7.2\% & 5.1\% \\
  \end{lstlisting}
  \begin{outputbox}
    \small
    \rowcolors{1}{gray!12}{white}
    \begin{tabular}{lcc}
      Bangladesh & 6.0\% & 5.4\% \\
      India & 7.2\% & 5.1\% \\
    \end{tabular}
  \end{outputbox}
  \whybox{Most journal templates already set page size and margins --
  avoid overriding them unless a submission guideline explicitly says
  to. Colour should be used sparingly and only where the journal
  allows it: shading alternating table rows (as above, via the
  \texttt{xcolor} \texttt{table} option) is one of the few colour uses
  widely accepted in print. Print-only journals often require
  greyscale -- when in doubt, the journal's ``Guide for Authors''
  overrides anything covered today.}
\end{frame}

\section{Migrating a Word Draft}

\begin{frame}{Converting an Existing Word Draft to LaTeX}
  \small
  \begin{enumerate}
    \item Start from the chosen journal template -- don't build from a blank file.
    \item Move text section by section (Introduction, Methodology, \dots), matching your Day 5 structure.
    \item Rebuild equations natively (Day 2) rather than pasting equation-editor images.
    \item Rebuild tables with \texttt{tabular}/\texttt{booktabs} (Day 3) -- do not paste Word tables as images.
    \item Re-attach your \texttt{.bib} file and citations (Day 4) last, once the text is in place.
  \end{enumerate}
\end{frame}

\section{Practice}

\begin{frame}{Hands-on Exercise}
  \begin{exercisebox}
    Choose one real economics journal template from the Overleaf
    gallery and migrate your \textbf{entire} running paper draft into
    it end-to-end: text, equations, tables, citations, and document
    structure all in one file (or split via \texttt{\textbackslash input} as in Day 5).
  \end{exercisebox}
\end{frame}

\begin{frame}{Homework Before Day 7}
  \begin{homeworkbox}
    Proofread the migrated document for formatting breaks: overflowing
    tables, misplaced figures, citation style mismatches, or JEL/keyword
    fields left as placeholder text. List everything you find -- you'll
    fix it tomorrow.
  \end{homeworkbox}
  \vspace{10pt}
  \small\textbf{Resources for today:}
  \begin{itemize}\small
    \item Overleaf Template Gallery -- \texttt{overleaf.com/latex/templates} (filter: academic-journal)
    \item Overleaf -- ``Page size and margins'' / ``Using colours in LaTeX''
    \item Wikibooks LaTeX -- Customizing chapters
  \end{itemize}
\end{frame}

\begin{frame}{Coming Up: Day 7}
  \Large
  \centering
  Troubleshooting, Polish \& Independent Mastery \\[6pt]
  \normalsize
  Reading errors like a pro, a final proofreading pass, and
  becoming self-sufficient beyond this course.
\end{frame}

\begin{frame}[plain]
  \begin{center}
    \vspace{35pt}
    {\Large\bfseries Thank You}\\[10pt]
    {\large From \texttt{\textbackslash begin} to \texttt{\textbackslash end}: The Complete LaTeX Guide}\\[4pt]
    Day 6 of 7 --- Journal Templates \& Professional Formatting\\[20pt]
    \rule{0.5\textwidth}{0.4pt}\\[10pt]
    Atikur Rahman \\
    White Globe Initiative \\[2pt]
    \texttt{atik@whiteglobeinitiative.org} \\
    \texttt{contact@whiteglobeinitiative.org}
  \end{center}
\end{frame}

\end{document}
